\subsection{Ordered Component Storage and Positional Access}

A sequence stores components in a fixed left-to-right order at any instant. The three built-in sequence types studied here differ in mutability and payload semantics:

\begin{itemize}
    \item \texttt{str}: immutable sequence of Unicode codepoints;
    \item \texttt{list}: mutable sequence of \texttt{PyObject*} references;
    \item \texttt{tuple}: immutable sequence of \texttt{PyObject*} references.
\end{itemize}

Syntax such as \texttt{len(obj)}, \texttt{obj[i]}, and \texttt{x in obj} delegates to special methods (\texttt{\_\_len\_\_}, \texttt{\_\_getitem\_\_}, \texttt{\_\_contains\_\_}) on the receiving object. The surface syntax is uniform; the C-level implementations are not.

\subsubsection{Slicing: Extracting Sub-Sequences with \texttt{[start:stop:step]}}

Slicing is not string cutting for beginners. It is a \emph{dynamic index projection}: the interpreter evaluates \texttt{start}, \texttt{stop}, and \texttt{step} (each defaulting when omitted), normalizes negative indices against sequence length, and walks the index range with the given stride. For \texttt{str} and \texttt{list}, the result is a \emph{new heap object} whose payload is copied from the selected components---immutable strings always allocate fresh character storage; lists allocate a new pointer array and copy the referenced slots. For \texttt{range} objects, by contrast, slicing returns another lightweight \texttt{range} descriptor over adjusted bounds without materializing elements.

The stride governs pointer arithmetic in the abstract sense: step \texttt{1} visits consecutive logical indices; step \texttt{2} skips every other slot; a negative step reverses traversal direction and may require \texttt{start} and \texttt{stop} to lie on opposite sides of the sequence. A step of zero is rejected at evaluation time (\texttt{ValueError}) because no valid traversal exists. Out-of-range slice bounds clamp silently---unlike scalar indexing, which raises \texttt{IndexError}---because the slice machinery treats bounds as half-open interval endpoints rather than single-address probes.

At the C level, \texttt{PySequence\_GetSlice} (and type-specific fast paths) computes output length from the normalized index triple, allocates a result object of the same sequence type, and copies either code units (\texttt{str}) or reference slots (\texttt{list}, \texttt{tuple}). The original object is never mutated. Slicing therefore always implies at least one allocation for non-empty results on materialized sequence types, and its cost scales with the number of extracted components, not with a constant-time pointer adjustment.
